\documentclass[12pt]{article}
\usepackage[T1]{fontenc}
\usepackage[utf8]{inputenc}
\usepackage{lmodern}
\usepackage{textcomp}
\usepackage{enumitem}
\usepackage{geometry}
\geometry{margin=1in}
\begin{document}
\begin{center}
Answer Key - Psychology - Graduate Level Assessment 9 \
\small (Answers are ordered exactly as in the question paper.)
\end{center}

\section*{Multiple Choice}
\begin{enumerate}[label=\arabic*.]
\item \textbf{Answer:} B
\\ \textit{Options:} A) left occipital; B) left temporal; C) both frontal lobes; D) right occipital; E) left parietal
\\ \textit{Note:} The left temporal lobe is crucial for language comprehension, as it houses Wernicke's area, which is responsible for the processing and understanding of spoken and written language.
\item \textbf{Answer:} A
\\ \textit{Options:} A) Extrinsic goal; B) Ultimate goal; C) Abstract goal; D) Intrinsic goal; E) Fundamental goal
\\ \textit{Note:} Money is considered an extrinsic goal because it is an external motivator that individuals pursue to achieve rewards or outcomes beyond the inherent satisfaction of the activity itself.
\item \textbf{Answer:} A
\\ \textit{Options:} A) Have no significant effect on the performance of tasks; B) Facilitate the performance of both easy and difficult tasks; C) hinder the performance of all tasks; D) Hinder the performance of difficult tasks; E) hinder the performance of easy tasks
\\ \textit{Note:} The correct answer reflects the concept of social facilitation and inhibition, where the presence of others does not significantly impact performance on well-learned or simple tasks, as extensive research demonstrates that audience effects vary depending on task complexity and individual familiarity.
\item \textbf{Answer:} A
\\ \textit{Options:} A) Mood swings, depression, anxiety; B) Social media usage, internet browsing habits, television viewing preferences; C) Intelligence, memory, reaction time; D) Speech patterns, tone of voice, cadence; E) Happiness, sadness, anger
\\ \textit{Note:} Mood swings, depression, and anxiety are all recognized indicators of emotional states, as they reflect significant fluctuations or disturbances in a person's emotional well-being and are commonly studied in psychological research.
\item \textbf{Answer:} E
\\ \textit{Options:} A) adults with speech impairments.; B) children under the age of two.; C) adults who are mentally retarded.; D) adults with physical disabilities.; E) children with hearing impairments.
\\ \textit{Note:} The Leiter International Performance Scale is designed to assess non-verbal intelligence, making it particularly useful for children with hearing impairments who may struggle with verbal communication.
\end{enumerate}

\section*{True / False}
\begin{enumerate}[label=\arabic*.]
\setcounter{enumi}{5}
\item \textbf{Answer:} False
\\ \textit{Options:} A) True; B) False
\\ \textit{Note:} The statement is false because the mode refers to the value that occurs most frequently in a dataset, and if the mode is stated as 18, it must be supported by the data showing that 18 appears more often than any other value, which is not confirmed in the statement.
\item \textbf{Answer:} True
\\ \textit{Options:} A) True; B) False
\\ \textit{Note:} A functional analysis is designed to systematically assess the relationship between specific environmental variables and behavior, thereby identifying the controlling variables that influence that behavior.
\item \textbf{Answer:} True
\\ \textit{Options:} A) True; B) False
\\ \textit{Note:} The behavioral differences observed in identical twins raised apart can largely be attributed to their distinct environments because environmental factors, such as upbringing, culture, and social interactions, significantly influence personality development and behavior, even in genetically identical individuals.
\item \textbf{Answer:} False
\\ \textit{Options:} A) True; B) False
\\ \textit{Note:} The heritability of traits in a cloned population is considered to be near 100\% because clones are genetically identical, which means that the variation in traits is primarily due to environmental factors rather than genetic differences.
\item \textbf{Answer:} True
\\ \textit{Options:} A) True; B) False
\\ \textit{Note:} Research indicates that individuals with higher levels of openness to experience are more imaginative, creative, and willing to engage in novel experiences, which correlates with a greater susceptibility to hypnosis.
\end{enumerate}

\section*{Long Form Response}
\begin{enumerate}[label=\arabic*.]
\setcounter{enumi}{10}
\item \textbf{Answer:} The parietal lobe plays a crucial role in integrating sensory information and is involved in the processing of language, which is essential for effective communication. Specifically, areas within the parietal lobe, such as the supramarginal gyrus, contribute to the synthesis of auditory and visual stimuli, facilitating speech production and comprehension. Alcohol consumption can impair the functioning of the parietal lobe, leading to difficulties in speech articulation and reduced clarity in verbal communication. This impairment is often manifested in slurred speech and a decreased ability to construct coherent sentences, highlighting the lobe's importance in both motor control and linguistic abilities. Consequently, understanding the impact of alcohol on the parietal lobe can provide insights into the broader effects of substance use on cognitive and communicative functions.
\\ \textit{Note:} correct because it accurately describes the parietal lobe's role in language processing and sensory integration, as well as how alcohol impairs its functions, leading to communication deficits.
\item \textbf{Answer:} Major injuries to the medulla oblongata can severely impair physiological functions, particularly swallowing, due to its critical role in regulating autonomic processes, including the swallowing reflex. Dysphagia, or difficulty swallowing, can arise from such injuries, leading to significant challenges in nutrition and hydration, which may necessitate modifications in diet or feeding methods. The psychological impact of these impairments can be profound; individuals may experience anxiety, depression, and social isolation due to the fear of choking or the embarrassment associated with eating difficulties. The resultant changes in daily life can disrupt social interactions and familial dynamics, as mealtimes are often central to social engagement. Thus, the interplay between physiological dysfunction and psychological well-being highlights the importance of comprehensive care for individuals affected by medullary injuries.
\\ \textit{Note:} correct because it accurately identifies the medulla oblongata's essential role in autonomic functions like swallowing, explains the resulting challenges such as dysphagia, and acknowledges the significant psychological effects on individuals' daily lives, including anxiety and depression.
\end{enumerate}

\vfill
\noindent\textit{End of Answer Key}
\end{document}